\documentclass[letterpaper]{article}

\usepackage[margin=1in]{geometry}
\usepackage{times}
\usepackage{amsmath}
\usepackage{url}

\urlstyle{rm}
\frenchspacing

\title{AAAI Reproducibility Checklist\\DriftingMol}
\author{Jiangjie Qiu, Yijun Li, Wentao Li, Xiaonan Wang\\
Beijing Key Laboratory of Artificial Intelligence for\\
Advanced Chemical Engineering Materials}
\date{}

\begin{document}
\maketitle

This checklist is stored separately from the main paper PDF.

\begin{enumerate}
    \item \textbf{Claims and scope.} The paper states that DriftingMol is a
    mechanism ablation for one-forward-pass property-biased SELFIES latent
    generation, not a molecular-generation state-of-the-art claim.
    \item \textbf{Theory.} Assumptions and proof sketches for decoder coupling,
    the induced pullback metric, z-diversity, and inherited SELFIES validity are
    given in the Mathematical Analysis section.
    \item \textbf{Data.} Experiments use ZINC250K with a fixed
    199,564/24,945/24,946 train/validation/test latent-cache split. The
    evaluated properties are QED, SA, LogP, and molecular weight.
    \item \textbf{Metrics.} The paper reports validity, uniqueness, novelty,
    Spearman correlation, MAE, calibration slope, internal diversity, scaffold
    diversity, and seed-level confidence intervals where available. RDKit parse
    failures count against validity.
    \item \textbf{Experimental protocol.} Conditional checkpoints are evaluated
    with 10,000 generated molecules over
    $\alpha\in\{1.0,1.5,2.0,3.0,5.0\}$, with the gate V$\ge$95\%,
    U$\ge$10\%, N$\ge$95\%. Matched retrieval and VAE-jitter references are
    evaluated with the same QED target bins and metric code.
    \item \textbf{Hyperparameters and compute.} Core model sizes, loss weights,
    temperatures, guidance grid, VAE reconstruction quality, and inference
    throughput are reported in the paper; full run configurations and generated
    result artifacts are maintained in the supplementary repository package.
    \item \textbf{Randomness.} The key QED settings use seeds 42--44: default
    DriftingMol, single-temperature decoder drift, no-diversity decoder drift,
    and the balanced multi-layer decoder setting. Other ablations and the fair
    four-property table are single seed unless otherwise stated.
    \item \textbf{Baselines.} The submission includes internal ablations,
    nonparametric matched references, and explicit limitations explaining the
    absence of protocol-matched external molecular-generator baselines.
    \item \textbf{Human subjects and societal impact.} No human-subject data are
    used. Molecular generation is dual-use; generated molecules require
    downstream filtering and expert review before experimental use.
\end{enumerate}

\end{document}
